%----------------------------------------------------------------------------------------
% Tailored CV — PhD Candidate, AMBIOM (Analysis of Microscopic BIOMedical images)
% Leibniz-Institut für Analytische Wissenschaften (ISAS), Dortmund.
% Ref 401_2026. Deadline 08.07.2026. TV-L E13, 4 years.
% Lead: microscopy image analysis / computer vision + ML representation learning + biology.
%----------------------------------------------------------------------------------------

\documentclass[10pt,a4paper,sans]{moderncv}

\usepackage[utf8]{inputenc}
\usepackage[T1]{fontenc}
\usepackage{lmodern}

\moderncvstyle{classic}
\moderncvcolor{blue}

\usepackage[scale=0.78]{geometry}
\usepackage{ragged2e}

%----------------------------------------------------------------------------------------

\firstname{Kundai Farai}
\familyname{Sachikonye}

\address{Auf der Lichtung 19}{82178 Puchheim, Germany}
\mobile{(+49) 1626386344}
\email{kundai.sachikonye@bitspark.com}
\homepage{github.com/fullscreen-triangle}
\photo[70pt][0.4pt]{pictures/Bew2.jpg}

\quote{Microscopy image analysis and computer vision, with machine-learning
representation learning for cell modelling --- and the biology to ground it.}

%----------------------------------------------------------------------------------------

\begin{document}

\makecvtitle

%----------------------------------------------------------------------------------------
%	PROFILE
%----------------------------------------------------------------------------------------

\section{Profile}

\cvitem{}{
  Computational researcher in \textbf{microscopy image analysis and computer
  vision}, with hands-on \textbf{machine-learning} and a background in
  \textbf{cell biology and bioinformatics}. I have built and deployed a
  quantitative microscopy image-analysis framework, and I am keen to formalise
  this work into a PhD on large-scale 3D microscopy analysis, shape modelling,
  and representation learning for cells.
}

%----------------------------------------------------------------------------------------
%	SELECTED WORK
%----------------------------------------------------------------------------------------

\section{Selected Work}

\cvitem{Microscopy image analysis}{
  \textbf{Quantitative image-analysis framework} --- segmentation, point-spread-function
  deconvolution, scale-field / coordinate (geometry) reconstruction, and
  morphometry with uncertainty quantification, for fluorescence and
  phase-contrast microscopy. GPU-accelerated (WebGL/GLSL); live, deployed tool.
  \href{https://hieronymus-omega.vercel.app/}{hieronymus-omega.vercel.app}
}

\cvitem{Shape \& geometry modelling}{
  \textbf{Coordinate-field reconstruction} --- recovering metric-preserving
  geometry from image observations (scale fields, geodesic distance), the basis
  for quantitative 3D shape characterisation of biological structures.
}

\cvitem{ML / representation learning}{
  Deep learning in \textbf{PyTorch} (CNNs, autoencoders / representation learning,
  attention); foundation-model fine-tuning (LoRA); applied to spectral and
  imaging data and to data-driven cell/disease-state modelling.
}

\cvitem{Cell modelling}{
  \textbf{Cell-state characterisation} --- spectral and circuit models of cellular
  state for diagnostics (tumour disease state, immunopeptidomics), connecting
  imaging and molecular data.
}

%----------------------------------------------------------------------------------------
%	EXPERIENCE
%----------------------------------------------------------------------------------------

\section{Experience}

\cventry{2021--present}{Independent Computational Researcher}{Bitspark GmbH}{}{}{
  Full-time research in microscopy image analysis, computer vision, machine
  learning, and computational biology. Built and deployed open, documented,
  GPU-accelerated tools; authored scientific papers.
}

\cventry{2019--2021}{Computational Lipidomics --- Bioinformatics \& HPC}{Technische Universität München}{Munich}{}{
  Large-scale high-throughput data analysis at HPC scale; CNN-based classification
  (PyTorch); statistical evaluation of $>$100\,GB datasets. Taught at Master's
  level and supervised student work.
}

\cventry{2017--2018}{Glycomics and Proteomics}{Universitätsklinikum Hamburg-Eppendorf}{Hamburg}{}{
  Automated high-throughput LC-MS/MS and MALDI-TOF data pipelines.
}

%----------------------------------------------------------------------------------------
%	EDUCATION
%----------------------------------------------------------------------------------------

\section{Education}

\cventry{2018--2019}{MSc Computational Biology}{Universität Tübingen}{}{}{
  Image and sequence analysis, statistics, and machine learning for biological
  data.
}
\cventry{2014--2017}{MSc Pharmaceutical Biotechnology}{HAW Hamburg}{}{}{
  Molecular and cell biology, analytical methods.
}
\cventry{2011--2014}{BSc Biochemistry and Cell Biology}{Jacobs University Bremen}{}{}{
  Thesis: DNA interstrand crosslinking --- cell culture and fluorescence
  microscopy.
}

\cvitem{}{\textit{Additional:} doctoral research in computational lipidomics at
  TUM (2019--2021; not completed).}

%----------------------------------------------------------------------------------------
%	TECHNICAL SKILLS
%----------------------------------------------------------------------------------------

\section{Technical Skills}

\cvitem{Image analysis / CV}{
  Segmentation, deconvolution, registration, morphometry; coordinate-field /
  geometry reconstruction; GPU/GLSL image processing; 2D and (towards) 3D
}

\cvitem{Machine learning}{
  PyTorch (CNNs, autoencoders, attention); representation learning;
  foundation-model fine-tuning (LoRA); scikit-learn
}

\cvitem{Programming}{
  Python (primary), Rust, TypeScript/WebGL, R, Bash, SQL
}

\cvitem{Biology}{
  Cell biology, biochemistry, bioinformatics; transcriptomics, proteomics,
  metabolomics
}

%----------------------------------------------------------------------------------------
%	LANGUAGES
%----------------------------------------------------------------------------------------

\section{Languages}

\cvitemwithcomment{English}{Native / Fluent (C2)}{}
\cvitemwithcomment{German}{Conversational}{resident in Germany since 2011}

%----------------------------------------------------------------------------------------

\renewcommand{\listitemsymbol}{-~}

\end{document}
